% GENERATED FILE. Edit canonical lessons, then run npm run generate:books.
% canonical-source-hash: fnv1a64:5bb336ab09f9490a
% canonical-lessons: IT-C492-abbraccio, IT-C492-coccolare, IT-C492-coccole, IT-C492-affezionarsi, IT-C492-affettuoso

\chapter{Abbracci e coccole}
\label{ch:it-abbracci-e-coccole}
\hlchaptermodality{492}

\begin{chapteropening}
\textbf{By the end of this chapter:} \emph{I can talk about affection: hugs, cuddles and growing fond of someone.}

Complete the last lesson of chapter 492: I can talk about affection: hugs, cuddles and growing fond of someone.
\end{chapteropening}

\section[l'abbraccio]{l'abbraccio — hug, embrace}

\label{lesson:IT-C492-abbraccio}



\begin{quote}
Before the new one: say the Italian for tenderness, then the Italian for kiss.
\end{quote}

\subsection*{You'll want to know: l'abbraccio}
\textbf{l'abbraccio} — \textquotedblleft{}hug, embrace\textquotedblright{}.

\begin{quote}
Mia nonna mi saluta sempre con un abbraccio forte.
\end{quote}

\begin{quote}
\emph{My grandmother always greets me with a big hug.}
\end{quote}

The noun for \emph{abbracciare}. At the end of a message: \emph{un abbraccio} — a hug.

\begin{cousinweb}
\emph{abbraccio} comes from \emph{abbracciare} \textquotedblleft{}to hug\textquotedblright{}, built on \emph{braccio} \textquotedblleft{}arm\textquotedblright{}, which is inherited from Latin \emph{bracchium} (Lewis \& Short).
\end{cousinweb}

\begin{grammarlens}[title={What you've built}]
One more word for telling, explaining and describing at B1.
\end{grammarlens}

\subsection*{Your turn}
\begin{itemize}
\raggedright
  \item \emph{Say it:} \emph{l'abbraccio}
  \item \emph{Say it:} \emph{l'abbraccio}, once more
  \item \emph{Say it:} a sentence of your own with \emph{l'abbraccio}
  \item \emph{Recall:} say the Italian for tenderness, then the Italian for kiss, then say \emph{l'abbraccio} again
\end{itemize}

\begin{tcolorbox}[breakable,colback=teal!4,colframe=teal!35!black,title={Before you move on}]
What does l'abbraccio mean? (\textquotedblleft{}Hug, embrace\textquotedblright{}.) Say it once more.
\end{tcolorbox}
\section[coccolare]{coccolare — to cuddle, to pamper}

\label{lesson:IT-C492-coccolare}



\begin{quote}
Before the new one: say the Italian for kiss, then the Italian for hug.
\end{quote}

\subsection*{You'll want to know: coccolare}
\textbf{coccolare} — \textquotedblleft{}to cuddle, to pamper\textquotedblright{}.

\begin{quote}
Il nonno coccola il nipotino tutto il giorno.
\end{quote}

\begin{quote}
\emph{Grandpa cuddles his little grandson all day.}
\end{quote}

Also to spoil with attention: \emph{coccolarsi} — to treat oneself.

\begin{cousinweb}
The origin of \emph{coccolare} is not settled; it may come from a word children use.
\end{cousinweb}

\begin{grammarlens}[title={What you've built}]
One more word for telling, explaining and describing at B1.
\end{grammarlens}

\subsection*{Your turn}
\begin{itemize}
\raggedright
  \item \emph{Say it:} \emph{coccolare}
  \item \emph{Say it:} \emph{coccolare}, once more
  \item \emph{Say it:} a sentence of your own with \emph{coccolare}
  \item \emph{Recall:} say the Italian for kiss, then the Italian for hug, then say \emph{coccolare} again
\end{itemize}

\begin{tcolorbox}[breakable,colback=teal!4,colframe=teal!35!black,title={Before you move on}]
What does coccolare mean? (\textquotedblleft{}To cuddle, to pamper\textquotedblright{}.) Say it once more.
\end{tcolorbox}
\section[le coccole]{le coccole — cuddles}

\label{lesson:IT-C492-coccole}



\begin{quote}
Before the new one: say the Italian for hug, then the Italian for to cuddle.
\end{quote}

\subsection*{You'll want to know: le coccole}
\textbf{le coccole} — \textquotedblleft{}cuddles\textquotedblright{}.

\begin{quote}
Il gatto vuole le coccole.
\end{quote}

\begin{quote}
\emph{The cat wants cuddles.}
\end{quote}

Usually plural: \emph{fare le coccole a qualcuno}.

\begin{cousinweb}
\emph{coccole} is the noun that goes with \emph{coccolare} \textquotedblleft{}to cuddle\textquotedblright{}; the origin of the word is not settled.
\end{cousinweb}

\begin{grammarlens}[title={What you've built}]
One more word for telling, explaining and describing at B1.
\end{grammarlens}

\subsection*{Your turn}
\begin{itemize}
\raggedright
  \item \emph{Say it:} \emph{le coccole}
  \item \emph{Say it:} \emph{le coccole}, once more
  \item \emph{Say it:} a sentence of your own with \emph{le coccole}
  \item \emph{Recall:} say the Italian for hug, then the Italian for to cuddle, then say \emph{le coccole} again
\end{itemize}

\begin{tcolorbox}[breakable,colback=teal!4,colframe=teal!35!black,title={Before you move on}]
What does le coccole mean? (\textquotedblleft{}Cuddles\textquotedblright{}.) Say it once more.
\end{tcolorbox}
\section[affezionarsi]{affezionarsi — to grow fond of, to get attached to}

\label{lesson:IT-C492-affezionarsi}



\begin{quote}
Before the new one: say the Italian for to cuddle, then the Italian for cuddles.
\end{quote}

\subsection*{You'll want to know: affezionarsi}
\textbf{affezionarsi} — \textquotedblleft{}to grow fond of, to get attached to\textquotedblright{}.

\begin{quote}
Mi sono affezionato al cane dei vicini.
\end{quote}

\begin{quote}
\emph{I have grown fond of the neighbours' dog.}
\end{quote}

\emph{Affezionarsi a} + person or thing. \emph{Affezionato} — fond, attached.

\begin{cousinweb}
\emph{affezionarsi} is built on \emph{affezione} \textquotedblleft{}affection\textquotedblright{}, a learned borrowing from Latin \emph{affectiō} \textquotedblleft{}a state of mind, a feeling\textquotedblright{} (Lewis \& Short), from \emph{afficere} \textquotedblleft{}to affect\textquotedblright{}.
\end{cousinweb}

\begin{grammarlens}[title={What you've built}]
One more word for telling, explaining and describing at B1.
\end{grammarlens}

\subsection*{Your turn}
\begin{itemize}
\raggedright
  \item \emph{Say it:} \emph{affezionarsi}
  \item \emph{Say it:} \emph{affezionarsi}, once more
  \item \emph{Say it:} a sentence of your own with \emph{affezionarsi}
  \item \emph{Recall:} say the Italian for to cuddle, then the Italian for cuddles, then say \emph{affezionarsi} again
\end{itemize}

\begin{tcolorbox}[breakable,colback=teal!4,colframe=teal!35!black,title={Before you move on}]
What does affezionarsi mean? (\textquotedblleft{}To grow fond of, to get attached to\textquotedblright{}.) Say it once more.
\end{tcolorbox}
\section[affettuoso]{affettuoso — affectionate, loving}

\label{lesson:IT-C492-affettuoso}



\begin{quote}
Before the new one: say the Italian for cuddles, then the Italian for to grow fond of.
\end{quote}

\subsection*{You'll want to know: affettuoso}
\textbf{affettuoso} — \textquotedblleft{}affectionate, loving\textquotedblright{}.

\begin{quote}
Mia figlia è molto affettuosa con tutti.
\end{quote}

\begin{quote}
\emph{My daughter is very affectionate with everyone.}
\end{quote}

At the end of a letter: \emph{saluti affettuosi} — warm wishes.

\begin{cousinweb}
\emph{affettuoso} is a learned borrowing from Late Latin \emph{affectuōsus} \textquotedblleft{}full of affection\textquotedblright{} (Lewis \& Short), from \emph{affectus} \textquotedblleft{}feeling\textquotedblright{}.
\end{cousinweb}

\begin{grammarlens}[title={What you've built}]
One more word for telling, explaining and describing at B1.
\end{grammarlens}

\subsection*{Your turn}
\begin{itemize}
\raggedright
  \item \emph{Say it:} \emph{affettuoso}
  \item \emph{Say it:} \emph{affettuoso}, once more
  \item \emph{Say it:} a sentence of your own with \emph{affettuoso}
  \item \emph{Recall:} say the Italian for cuddles, then the Italian for to grow fond of, then say \emph{affettuoso} again
\end{itemize}

\begin{tcolorbox}[breakable,colback=teal!4,colframe=teal!35!black,title={Before you move on}]
What does affettuoso mean? (\textquotedblleft{}Affectionate, loving\textquotedblright{}.) Say it once more.
\end{tcolorbox}
